% from ega3-1-fr.tex
\oldpage[III]{136}
on peut supposer que $f^{-1}(U)=V$, d'où on conclut aussitôt, par définition,
que $V\subset X'$, ce qui achève la démonstration.
\end{proof}

La proposition suivante a été démontrée par Chevalley dans le cas des schémas
algébriques~:

\begin{proposition}[4.4.2]
\label{III.4.4.2-fr}
Soient $Y$ un préschéma localement noethérien, $f:X\to Y$ un morphisme. Les
conditions suivantes sont équivalentes~:
\begin{enumerate}
\item[\textup{a)}] $f$ est fini.
\item[\textup{b)}] $f$ est affine et propre.
\item[\textup{c)}] $f$ est propre et, pour tout $y\in Y$, $f^{-1}(y)$ est un ensemble
fini.
\end{enumerate}
\end{proposition}

\begin{proof}
On sait que \emph{a)} entraîne \emph{b)} (\textbf{II}, 6.1.2 et 6.1.11). Si $f$ est propre et affine,
il en est de même du morphisme
$f^{-1}(y)\to\operatorname{Spec}(\boldsymbol{k}(y))$ (\textbf{II}, 1.6.2, (iii) et 5.4.2, (iii)),
et le th. de finitude (\hyperref[III.3.2.1-fr]{3.2.1}) appliqué au faisceau
structural de $f^{-1}(y)$, montre que $f^{-1}(y)=\operatorname{Spec}(A)$, où
$A$ est une $\boldsymbol{k}(y)$-algèbre \emph{finie}~; donc $f^{-1}(y)$ est un ensemble fini
(\textbf{II}, 6.1.7), et on voit que \emph{b)} entraîne \emph{c)}. Enfin, comme $f^{-1}(y)$ est un
préschéma algébrique sur $\boldsymbol{k}(y)$, l'hypothèse que l'ensemble $f^{-1}(y)$ est
fini entraîne que l'espace $f^{-1}(y)$ est discret (\textbf{I}, 6.4.4). Avec les
notations de (\hyperref[III.4.4.1-fr]{4.4.1}), on a donc $X'=X$, et
$f':X\to Y'$ est un isomorphisme~; comme $g$ est un morphisme fini, on voit
que \emph{c)} entraîne \emph{a)}.
\end{proof}

\begin{theorem}[4.4.3]
\label{III.4.4.3-fr}
\textnormal{(«~Main theorem~» de Zariski).}
Soient $Y$ un préschéma noethérien, $f:X\to Y$ un morphisme quasi-projectif,
$X'$ l'ensemble des points $x\in X$ qui sont isolés dans leur fibre
$f^{-1}(f(x))$. Alors $X'$ est une partie ouverte de $X$, et le
sous-préschéma induit $X'$ est isomorphe à un préschéma induit sur une partie
ouverte d'un $Y$-préschéma $Y'$ fini sur $Y$.
\end{theorem}

\begin{proof}
L'hypothèse entraîne qu'il existe un $Y$-préschéma projectif $Z$ tel que $X$
soit $Y$-isomorphe à un sous-préschéma induit sur un ouvert de $Z$
(\textbf{II}, 5.3.2 et 5.5.1). On est donc ramené à démontrer le théorème lorsque $f$
est un morphisme projectif, donc propre (\textbf{II}, 5.5.3), et il résulte alors
aussitôt de (\hyperref[III.4.4.1-fr]{4.4.1}).
\end{proof}

\begin{remark}[4.4.4]
\label{III.4.4.4-fr}
Si $X$ est réduit (resp. irréductible et $X'$ non vide), on peut supposer,
dans l'énoncé de (\hyperref[III.4.4.3-fr]{4.4.3}), que $Y'$ est réduit (resp.
irréductible). En effet, on peut toujours remplacer $Y'$ par le
sous-préschéma \emph{adhérence} $\overline{X'}$ de $X'$ dans $Y'$ (\textbf{I}, 9.5.11 et
\textbf{II}, 6.1.5, (i) et (ii)), et on sait que si $X'$ est réduit, il en est de même
de $\overline{X}'$ (\textbf{I}, 9.5.9, (i))~; par ailleurs, si $X'$ n'est pas vide, il
est irréductible si $X$ l'est, et $\overline{X'}$ est alors aussi
irréductible.
\end{remark}

\begin{corollary}[4.4.5]
\label{III.4.4.5-fr}
Soient $Y$ un schéma localement noethérien, $f:X\to Y$ un morphisme de type
fini, $x$ un point de $X$ isolé dans sa fibre $f^{-1}(f(x))$. Alors il existe
un voisinage ouvert de $x$ dans $X$ qui est isomorphe à une partie ouverte
d'un $Y$-préschéma fini sur $Y$.
\end{corollary}

\begin{proof}
Soient en effet $y=f(x)$, $U$ un voisinage ouvert affine de $y$ dans $Y$,
$V$ un voisinage ouvert affine de $x$ dans $X$, contenu dans $f^{-1}(U)$.
Comme $Y$ est séparé, l'injection $U\to Y$ est affine (\textbf{II}, 1.6.3), et comme
$V$ est affine sur $U$ (\emph{ibid.}), la restriction de $f$ à $V$ est un
morphisme affine $V\to Y$ (\textbf{II}, 1.6.2, (ii))~; \emph{a fortiori}, cette
restriction est un morphisme quasi-projectif puisqu'il est de type fini
(\textbf{I}, 6.3.5 et \textbf{II}, 5.3.4, (i)). Il suffit alors d'appliquer à cette restriction
le th. (\hyperref[III.4.4.3-fr]{4.4.3}).
\end{proof}

Le cor. (\hyperref[III.4.4.5-fr]{4.4.5}) s'énonce dans le langage de
l'algèbre commutative~:

